\section*{Цель работы}
Экспериментальное исследование линейных разветвленных резистивных цепей с использованием методов наложения, эквивалентного источника и принципа взаимности.

\section*{Краткие теоретические сведения}
В работе анализируется резистивная цепь с источниками постоянного напряжения $U$ и тока $I$.

\begin{figure}[H]
    \centering
    \includegraphics[width=0.6\textwidth]{figs/circuit_basic.pdf}
    \caption{Исследуемая линейная резистивная цепь}
    \label{fig:circuit_basic}
\end{figure}

\textit{Метод наложения.} Реакция цепи на действие нескольких источников определяется как алгебраическая сумма реакций на действие каждого источника в отдельности.

\begin{figure}[H]
    \centering
    \includegraphics[width=0.8\textwidth]{figs/circuit_superposition.pdf}
    \caption{Пояснение метода наложения}
    \label{fig:circuit_superposition}
\end{figure}

\textit{Метод эквивалентного источника напряжения.} По отношению к одной из ветвей линейную цепь с несколькими источниками можно представить одним эквивалентным источником напряжения $U_0$ с последовательно соединенным сопротивлением $R_0$. Ток в рассматриваемой ветви определяется по закону Ома:
\begin{equation*}
    I_3 = \frac{U_0}{R_0 + R_3}
\end{equation*}

\begin{figure}[H]
    \centering
    \includegraphics[width=0.8\textwidth]{figs/circuit_thevenin.pdf}
    \caption{Схемы замещения для метода эквивалентного источника}
    \label{fig:circuit_thevenin}
\end{figure}

\textit{Принцип взаимности.} Если источник напряжения, действуя в одной ветви линейной цепи, вызывает ток в другой ветви, то тот же источник после переноса во вторую ветвь вызовет в первой ветви такой же ток.

\begin{figure}[H]
    \centering
    \includegraphics[width=0.8\textwidth]{figs/circuit_reciprocity.pdf}
    \caption{Иллюстрация принципа взаимности}
    \label{fig:circuit_reciprocity}
\end{figure}

\section*{Порядок выполнения работы (Снятие данных)}

Для выполнения работы используется экспериментальная плата, схема которой представлена на \cref{fig:circuit_experimental}.

\begin{figure}[H]
    \centering
    \includegraphics[width=0.7\textwidth]{figs/circuit_experimental.pdf}
    \caption{Экспериментальная схема установки}
    \label{fig:circuit_experimental}
\end{figure}

\subsection*{Опыт 1. Исследование цепи при питании от двух источников}
\begin{itemize}
    \item \textit{Что сделать:} Собрать схему и измерить все токи и напряжения ветвей.
    \item \textit{Как сделать:} 
    1. Подключить источник напряжения $U$ к гнездам 1 и 2 платы (плюс к 1, минус к 2).
    2. Включить источник напряжения $U$ переключателем S1 (положение вверх) и источник тока $I$ переключателем S2 (положение вверх).
    3. Установить $U = 4$ В (или 2 В по указанию).
    4. Измерить ток $I$ амперметром, включенным последовательно с источником тока.
    5. Измерить токи $I_1 \dots I_4$ и напряжения $U_1 \dots U_4$ во всех ветвях. Данные занести в \cref{tab:experiment_full}.
    \item \textit{Нюансы/Предупреждения:} Погрешность цифровых приборов меньше при фиксации положительных показаний. При отрицательных значениях меняйте полярность щупов и фиксируйте знак в таблице.
\end{itemize}

\subsection*{Опыт 2. Определение токов цепи методом наложения}
\begin{itemize}
    \item \textit{Что сделать:} Измерить частичные токи от каждого источника по отдельности.
    \item \textit{Как сделать:}
    1. Оставить включенным только источник $U = 4$ В (источник тока отключить переключателем S2). Измерить токи во всех ветвях.
    2. Оставить включенным только источник тока $I$ (источник $U$ отключить переключателем S1). Измерить токи во всех ветвях.
    3. Занести результаты в \cref{tab:superposition}.
    \item \textit{Нюансы/Предупреждения:} Обращайте внимание на направления токов. При отключении источника напряжения на плате происходит его закорачивание, а при отключении источника тока — разрыв цепи.
\end{itemize}

\subsection*{Опыт 3. Определение тока методом эквивалентного источника напряжения}
\begin{itemize}
    \item \textit{Что сделать:} Найти напряжение холостого хода $U_0$ и измерить ток $I_3$ при питании от эквивалентного источника.
    \item \textit{Как сделать:}
    1. Подключить оба источника ($U$ и $I$).
    2. Оборвать ветвь 3 на участке A-B. Измерить напряжение $U_0$ на разомкнутых зажимах A и B.
    3. Исключить из схемы источники $U$ и $I$ (переключатели S1 и S2 вниз).
    4. К разомкнутым гнездам A и B ветви 3 подключить источник напряжения и установить на нем значение $U = U_0$.
    5. Измерить ток $I_3$ и записать в протокол.
    \item \textit{Нюансы/Предупреждения:} При подключении источника $U_0$ к зажимам A-B строго соблюдайте полярность, определенную при измерении напряжения холостого хода.
\end{itemize}

\subsection*{Опыт 4. Экспериментальная проверка принципа взаимности}
\begin{itemize}
    \item \textit{Что сделать:} Убедиться в равенстве токов при взаимном обмене источника и амперметра.
    \item \textit{Как сделать:}
    1. Отключить источник тока $I$ от цепи.
    2. Установить заданное напряжение $U$ на входе и измерить ток $I_3$ в ветви 3.
    3. Отключить источник $U$ от входа и подключить его в разрыв ветви 3 (к гнездам A, B).
    4. Измерить ток $I_1$ в ветви 1. Записать данные в протокол.
    \item \textit{Нюансы/Предупреждения:} Значение напряжения переносимого источника $U$ должно быть абсолютно одинаковым в обоих измерениях.
\end{itemize}

% Соответствие изображений оригинальной методичке:
% fig:circuit_basic.pdf -> Рис. 2.1
% fig:circuit_superposition.pdf -> Рис. 2.2
% fig:circuit_thevenin.pdf -> Рис. 2.3
% fig:circuit_reciprocity.pdf -> Рис. 2.4
% fig:circuit_experimental.pdf -> Рис. 2.5